\section*{LỜI NÓI ĐẦU}
\thispagestyle{empty}
Ngày nay, khi công nghệ hiện đại phát triển nhanh chóng và được sử dụng trong mọi khía cạnh của cuộc sống, mọi người đều phải thích nghi với những gì đang thay đổi và tiếp thu cũng như phát triển những kỹ năng sâu rộng, đặc biệt là những người đang làm việc trong lĩnh vực công nghệ. Là sinh viên tại Đại học Bách Khoa Hà Nội, tôi phải học tập, nghiên cứu và sử dụng thành tựu công nghệ một cách hiệu quả. Từ lâu, các thiết bị tự hành đã giúp đỡ con người trong rất nhiều lĩnh vực từ sản xuất công nghiệp, nghiên cứu, cho đến đời sống thường nhật hằng ngày. Trong những năm gần đây, trí tuệ nhân tạo (AI) đã nhận được rất nhiều sự quan tâm của giới truyền thông, nhiều bài báo thường xuyên thảo luận về học máy, học sâu và AI. Chính vì vậy tôi quyết định chọn đề tài "Phân chia tính toán động với điện toán phân tán mã hóa và phân bổ năng lượng cho xe tự hành trang bị radar chức năng kép" cho nghiên cứu đồ án của tôi. Đây là kết quả của tôi trong toàn bộ học kỳ này, tôi đặt rất nhiều công sức cũng như thời gian cho đồ án này. Thông qua đồ án, tôi cung cấp một cái nhìn tổng quan toàn diện về hệ thống phương tiện tự hành và ứng dụng của học sâu nhằn cải thiện hiệu năng cũng như thời gian xử lý trong hệ thống, cũng như ý tưởng và cách tiếp cận đã được phát triển. 

Trường Đại học Bách Khoa Hà Nội đóng vai trò quan trọng trong cuộc đời tôi. Việc theo học tại đây đã mang lại cho tôi nhiều cơ hội phát triển, gặp gỡ các thầy cô và bạn bè tuyệt vời. Bên cạnh kiến thức và những cơ hội tương lai mở ra, Bách Khoa còn trang bị cho tôi những bài học quý giá trong quá trình học tập. Trong quá trình nghiên cứu, tôi đã tham gia vào Lab nghiên cứu của PGS.TS Nguyễn Tiến Hòa. Tại đây, tôi không chỉ học hỏi kiến thức lý thuyết và mô hình hệ thống phức tạp, mà còn nhận được những bài học quan trọng trong cuộc sống. Tôi xin chân thành cảm ơn PGS.TS Nguyễn Tiến Hòa và các thành viên trong Lab nghiên cứu đã đồng hành và hỗ trợ tôi trong quá trình nghiên cứu. Qua quá trình này, tôi đã tiếp cận với các mô hình viễn thông, trong đó có phương tiện tự hành. Thấy được tính ứng dụng cao và tầm quan trọng của hệ thống này, cùng với sự tò mò và đam mê với Trí tuệ nhân tạo và học máy, tôi đã quyết định theo đuổi và kết quả là đồ án này. Tôi muốn bày tỏ lòng biết ơn đến các cố vấn đã hướng dẫn và hỗ trợ tôi trong suốt quá trình nghiên cứu.

Tôi cũng muốn gửi lời cảm ơn đến gia đình và bạn bè vì sự động viên và sự hiểu biết của họ trong suốt những thử thách tôi đã trải qua.
% Nowadays, as modern technology is rapidly evolved and used in every aspects of life, everybody has to adapt to what is changing and be provided with extensive skills, especially the ones who are working in the technology field. As a student in HA NOI University of Science and Technology, I have to learn, understand how to use the world's technological achievements effectively. During the last several years, Artificial intelligence (AI) has received a lot of media attention, numerous articles frequently discuss machine learning, deep learning and AI. Due to this, I decide to choose the topic of Personalized Federated Learning for my thesis research. This is my culmination over the entire semester, I put a lot of hard work to complete  this. Through out the thesis, I provide a comprehensive overview of the current state-of-the-art in Federated Learning, Personalized Federated Learning, as well as the idea and approaches that been developed.

% Studying in Ha Noi University of Science and Technology is an important milestones in my life. I have chances to meet with many great people, many wonderful friends who will come alongside with me to the bright future. I joined signal processing lab of Prof Nguyen Tien Hoa. From this moment, I have gained alot of knowledge not only about the academic theory, complicated system model but also the lessons about live. I would like to express my sincere thanks Assoc. Prof. Nguyen Tien Hoa and all of my fellow lab members have helped to come to this journey. During the time working in the lab, I have experienced with machine learning and some interested form my side grown. I interact with Federated Learning and decide to pursue this concept as my thesis. 

% I would like to express my gratitude to my advisors and collaborators for their guidance and support throughout my research. I would also like to thank my family and friends for their encouragement and understanding during this challenging road. 
\cleardoublepage